\chapter{Discussion}
\label{chap:discussion}

The experiments show that graph-based vehicle-block construction is the most
important step in the implemented solution approach.  The greedy method already
finds feasible schedules, but it makes local decisions.  The path-cover methods
use the compatibility graph more systematically, and the multi-start weighted
version gives the best validated results among the implemented methods.  During
the experiments, this was one of the clearest observations: improving the way
trips are chained usually mattered more than changing only the fleet assignment
afterwards.

\section{Interpretation of the Main Findings}
\label{sec:discussion-main-findings}

The final archive contains feasible solutions for all twelve Senior instances.
The headline instances use \(2\), \(5\), and \(15\) vehicles for Small, Medium,
and Large, respectively.  Over the complete Senior benchmark, the final archive
has total validated cost \(10000.48\) and uses \(126\) vehicles.

The detailed tables and plots in \Cref{chap:results} report the behaviour of
the algorithms instance by instance.  For the discussion, the main outcome can
also be read in one compact view.  \Cref{fig:final_performance_dashboard}
combines the headline costs, the all-instance archive result, the fleet split,
the objective-audit consistency, and the computational setting.

\begin{figure}[H]
    \centering
    \includegraphics[width=\textwidth]{figures/final_performance_dashboard.pdf}
    \caption{Final result summary of the implemented multi-start path-cover matheuristic.}
    \label{fig:final_performance_dashboard}
\end{figure}

The result summary shows that the final method is strong because several indicators
point in the same direction.  It gives low validated costs on Small, Medium,
and Large, covers all twelve Senior instances, keeps the final archive at
\(10000.48\) cost with \(126\) used vehicle blocks, and does so in a CPU-only
local environment without GPU or HPC support.  The EV/ICE split also shows that
the method does not force electric operation everywhere.  EVs are used where
the block structure and charging opportunities make them feasible.

The results also show that electrification is not simply a matter of replacing
as many ICE vehicles as possible.  Medium can be operated with a 100\% EV share
in the final solution, while Small is cheaper with two ICE vehicles.  This is
not a contradiction.  In practice, EV use depends on the structure of the
blocks, available charging time, terminal capacity, and the fixed and
operational cost coefficients of the instance.

Fixed vehicle cost is the largest component in the final objective audit in
\Cref{tab:cost-audit-main}.  This explains why reducing the number of blocks is
often more important than small changes in deadhead or charging time.  At the
same time, charging feasibility cannot be ignored.  A low vehicle count is only
useful if the resulting blocks can still satisfy autonomy and terminal-capacity
requirements.
Initially, I expected the EV assignment step to have a larger influence on the
final objective.  The experiments made this point clearer: in most cases tested
here, the structure of the vehicle blocks had to be strong before EV assignment
could improve the solution.

The results therefore support a two-level interpretation.  At the first level,
the path-cover construction improves the schedule by joining compatible trips
and reducing unnecessary block starts.  This is visible in the vehicle counts
and in the dominance of fixed vehicle cost.  At the second level, the EV and
charging checks decide which of these blocks can actually be operated by
electric buses.  This second level does not always reduce the number of
vehicles, but it changes the fleet mix and the smaller cost components.

The Medium instance illustrates the favorable case for electrification.  The
final solution uses five vehicles and all of them are electric.  This means that
the block structure, break locations, charging times, and terminal capacities
fit the EV restrictions well.  The Small instance illustrates a different case.
The best validated solution uses only two vehicles, both conventional.  Adding
EV operation would not automatically improve the official objective if it
requires more vehicles or additional operational time.  The Large instance lies
between these two cases and uses a mixed fleet.

The all-instance results show the same pattern more broadly.  EV share,
vehicle productivity, and cost per trip are related but not identical.  A high
EV share can be operationally attractive, but it is not the only measure of
solution quality.  A schedule with fewer vehicles and lower cost can be better
under the official objective even if its EV share is lower.  This is why the
thesis reports EV share together with cost, vehicle count, charging minutes,
and objective components.
This may not hold in general for every operator or cost setting.  At least for
this benchmark, however, the result suggests that electrification has to be
discussed together with vehicle count and charging feasibility, not as a single
percentage.

\section{Benefits and Drawbacks of the Implemented Methods}
\label{sec:discussion-methods}

\Cref{tab:exact-heuristic-components} summarizes the methodological role of
the four implemented approaches.  The greedy method is useful as a transparent
baseline and as a quick feasibility constructor.  Its drawback is that early
successor choices can block better later combinations.

The unweighted path-cover method improves the search by solving a matching
problem on the complete compatibility graph.  It has a clear graph-theoretic
interpretation: more selected arcs usually means fewer vehicle blocks.  However,
it does not distinguish between a short clean continuation and a long
operationally unattractive continuation.

The weighted path-cover method addresses this weakness by scoring arcs with
waiting time, deadhead time, and bounded charging-risk information.  The score
is still a construction score and not the official objective.  Its advantage is
that the matching step becomes closer to the final cost calculation.  The
multi-start version adds diversification by testing several timetable and graph
variants.  In practice, this turned out to be helpful but not free: it makes the
search more stable, but it also increases runtime, especially for Large.

The comparison also shows why a single algorithmic idea is not enough for the
whole problem.  Greedy construction gives feasibility quickly, but it does not
look far enough ahead.  Matching improves the global structure of vehicle
blocks, but unweighted matching can choose unattractive continuations.  Weighted
matching improves the arc choice, but it still depends on the selected
timetable.  Multi-start search addresses this by changing the timetable and
graph variants, but it cannot explore all possible combinations.

The final method is therefore a compromise.  It uses exact graph optimization
where the structure is clear and limited, and it uses heuristic search where
the complete integrated problem would be too large.  This compromise is
appropriate for the MINOA Senior benchmark because the task contains several
interacting decisions: trip selection, block construction, fleet assignment,
battery propagation, charging insertion, and node-capacity feasibility.

The no-regression archive is another practical design choice.  It prevents a
new heuristic variant from replacing a better existing output.  This matters
because a heuristic can improve one instance and worsen another.  By retaining
the better validated output for each instance, the final archive reports the
best solution found by the implemented search process without claiming that one
single parameter setting is best for every instance.
This was not always the case during testing.  Some candidate variants improved
one file but produced a worse block structure on another, which is why the
archive rule became a useful reporting device rather than only an
implementation detail.

\section{Validation, Bounds, and Optimality}
\label{sec:discussion-validation-bounds}

The reported schedules are feasible upper bounds.  They are not claimed to be
globally optimal.  Matching is exact only for a fixed selected timetable,
fixed compatibility graph, and fixed arc weights.  The complete problem still
contains heuristic decisions in timetable generation, candidate-family
construction, EV assignment, and charging insertion.

The lower-bound results in \Cref{tab:headline-bound-gap} give useful context.
For Small, the global lower bound is close enough to show that the remaining
gap is moderate.  For Large, the controlled run does not provide a useful
certified global bound.  The selected-timetable lower bound is therefore only a
diagnostic value for the timetable already chosen by the heuristic.

The objective audit in \Cref{eq:validator-residual} is important for reporting
reliability.  It confirms that the internally computed cost and the external
audit value agree within rounding tolerance.  This does not prove optimality,
but it supports the claim that the reported objective values are computed
consistently.

The lower-bound results must be read carefully.  A feasible solution gives an
upper bound because it is an admissible schedule with a known cost.  A lower
bound gives a value that no feasible solution can beat for the corresponding
relaxation.  If both bounds coincide, global optimality is certified.  In this
thesis, the bounds do not close the complete integrated problem.  They are used
to interpret solution quality and to show where the remaining uncertainty is
larger.

The distinction between \(LB^{global}\) and \(LB^{TT}(\pi)\) is especially
important.  The global lower bound is valid before the timetable is fixed.  The
selected-timetable diagnostic lower bound assumes the timetable produced by the
heuristic.  It can therefore be much more informative for the vehicle-scheduling
subproblem, but it is not a certificate for the full integrated problem.  This
distinction avoids overstating the quality of the final schedules.

The MINOA desktop validator has a different role from the lower-bound calculation.
It checks whether a generated output satisfies the required format and
feasibility rules and whether the reported objective value is consistent.  It
does not certify that the solution is globally optimal.  This separation is
important for the scientific interpretation of the experiments: validation
confirms feasibility and objective agreement, while bounds provide information
about possible optimality gaps.

\section{Limitations and Future Work}
\label{sec:discussion-limitations}

The main limitation is the absence of a global optimality proof for the full
integrated problem.  The method searches over a controlled set of timetable and
graph variants, not over every possible integrated schedule.  A second
limitation is that charging-aware information enters the matching step through
bounded scoring and later feasibility checks rather than through one fully
integrated exact charging model.

Another limitation is that the multi-start family is bounded by design.  The
algorithm explores several timetable starts, graph scores, EV assignment
strategies, and repair options, but it does not perform an unrestricted search
over all successor swaps or all block recombinations.  This makes the method
reproducible and practical on the reported hardware, but it also means that a
better solution may exist outside the explored candidate family.

The charging model is also intentionally simplified compared with a detailed
electric-bus energy model.  The implementation checks autonomy, minimum
charging duration, slow and fast charging options, and terminal capacity.  It
does not model nonlinear charging curves, temperature effects, battery
degradation, or stochastic energy consumption.  These details are outside the
MINOA Senior input data used in this thesis, but they would matter in a
real-world deployment study.

The results are based on the available Senior instances.  They provide useful
evidence for the implemented methods, but they are not a statistical sample of
all possible bus networks.  Different cities may have other terminal layouts,
different depot locations, different charger capacities, and different fleet
cost coefficients.  The graph-based method should transfer conceptually, but
the numerical performance would need to be tested on those data sets.
When running the larger instances, this dependence on the input structure was
easy to see.  Small changes in line structure or charging opportunity can
change which blocks are attractive for EV operation.

\paragraph{Exact optimization and global optimality certification.}
\label{sec:future-global-optimality}

The proposed matheuristic provides feasible solutions and therefore valid upper
bounds.  A future exact method could use these solutions as warm starts or
initial upper bounds.  Such a method would need to model timetable selection,
vehicle connections, vehicle types, battery states, charging activities, and
parking and charging capacities jointly.  It should also avoid restricting the
search only to the timetable variants generated by the current heuristic.

Possible exact directions include mixed-integer programming, branch-and-cut,
branch-and-price, or decomposition.  In all cases, valid lower bounds would be
required in addition to feasible upper bounds.  If \(LB\) is a valid lower bound,
\(UB\) is the cost of the best feasible solution, and \(C^\star\) is the unknown
global optimum, then
\begin{equation}
  LB \le C^\star \le UB .
  \label{eq:future-bound-relation}
\end{equation}
Global optimality would be certified only when \(LB=UB\).  More commonly, an
exact solver may stop with a certified remaining optimality gap.  This extension
could require substantially more computation time for larger instances.

Future work could also combine the current matheuristic with an exact method.
The heuristic solution could provide a warm start and a strong initial upper
bound.  The exact method could then focus on proving whether a better solution
exists.  This would preserve the practical advantage of the current solver
while adding a stronger certification layer.

A second extension would be a richer charging-aware neighborhood search.  The
current implementation uses bounded local improvements and charging-aware
scores.  A larger neighborhood could simultaneously move trips between blocks,
change charging locations, and reassign EVs.  Such a neighborhood may improve
cost, but it would also require careful feasibility checks to avoid violating
capacity constraints.

A third extension would be multi-objective analysis.  The official objective
combines fixed vehicle cost, paid break cost, pull-in/out cost, and CO\(_2\)
cost.  Public transport operators may also care about separate indicators such
as EV share, charging-infrastructure utilization, operational stability, or
reserve battery level.  Reporting these indicators separately would make the
trade-offs between cost and sustainability more explicit.
